\section{Discussion}
\label{sec:discussion}

\subsection{Why stop-level lags work}

The stop-level lag features succeed because they capture real temporal
autocorrelation in the passenger count series without the leakage artifact
introduced by 1~Hz lags on forward-filled data. Each
$\text{stoplag}_k(t)$ value is the passenger count at a strictly earlier
stop event ($k$ stops ago), and the forward-fill broadcasting applied to
the \emph{lag values} (rather than the raw series before feature
computation) simply propagates this known earlier value through the
inter-stop interval. There is no pathway by which the current passenger
count can be recovered from the lag features, and Section~\ref{sec:results}
empirically confirms the mechanism carries real predictive signal (+0.22
macro-F1 across all three models). The PACF analysis (Figure~\ref{fig:pacf})
supports this by identifying lag~1 as the dominant direct temporal
dependency, consistent with the observation that lag-1 alone contributes
most of the Gini importance among the three stop-level lag features.

\subsection{Why TimesFM embeddings don't help}

The most striking observation is the 100\% PCA explained variance at 32
components. The foundation model's representations for this domain have at
most 32 independent dimensions of variation, out of the 1280 available in
the raw embedding space. Several interpretations are consistent with this
observation.

\textbf{Low-entropy domain.} Urban bus ridership is structured. There are
only so many ``shapes'' a bus trip can take: morning rush hour, evening
rush hour, weekday midday, weekend, night service, special events,
holidays. Each presents a distinctive temporal pattern, but the number of
distinct patterns is much smaller than 1280. The foundation model has been
pretrained on massive and diverse time series data, but at inference time
it projects our input into a subspace determined by what it has learned
\emph{about this kind of data}. For low-entropy domains, that subspace is
small.

\textbf{Redundancy with hand-crafted lags.} The stop-level lag features
already capture the most predictive component of the temporal structure,
the short-range autocorrelation at lag~1. Whatever additional information
the TimesFM embeddings encode (longer-range patterns, cross-mission
similarities, daily periodicity captured via attention over longer
contexts) is either not relevant to this specific classification task or
not linearly recoverable by tree ensemble models. Recent large-scale
benchmarks show that tree-based models consistently outperform neural
networks on tabular data~\cite{grinsztajn2022trees}, so the lack of
improvement from embedding augmentation is consistent with this broader
finding.

\textbf{Extraction lossiness.} Our double mean-pooling, first across the
sequence dimension of each context window, then across all windows of a
mission, is lossy. A more sophisticated aggregation (attention pooling,
learned temporal query, or per-stop embeddings rather than per-mission)
might preserve more signal. However, the simplicity of our pipeline is
itself informative: the ``drop in and train'' use case of foundation models
as feature extractors has a harder time with structured domains than with
less-structured ones like natural language, consistent with broader evidence that pretrained model representations
offer limited gains on structured time series
tasks~\cite{zheng2025llmtimeseries}. The null result is a useful
data point for practitioners considering whether to invest in foundation
model extraction pipelines for tabular time series tasks.

\textbf{Interpretability gap.} A natural diagnostic would be to
correlate each of the 32 retained PCA components against our
hand-crafted features to identify what the foundation model is
``seeing.'' We do not conduct this analysis, and we note that such a
correlation would be difficult to interpret even if computed: PCA
components on the 1280-dimensional transformer activation space are
orthogonal linear combinations of latent features that themselves lack
direct semantic grounding in the input time series. Our hand-crafted
features (e.g., \texttt{hour}, \texttt{pax\_stop\_lag\_1}) carry
explicit domain semantics; the PCA components do not. The 100\%
explained-variance-at-32-components result is the interpretable summary
we can offer: the foundation model's representations for this domain
live in a low-rank subspace, irrespective of what axes that subspace
happens to be oriented along.

\textbf{Design choice: PCA-reduced embeddings versus the forecast head.}
A valid alternative to PCA reduction would be to feed TimesFM's scalar
point forecast directly into the downstream classifier, since the
forecast head is itself a learned mapping from the 1280-dimensional
penultimate activations to a scalar prediction. This appears
superficially equivalent to PCA-then-RF (both compose compression with
a nonlinear head), but three considerations motivate our choice. First,
the forecast head was trained to minimize forecast error on future
passenger-count values, not to separate our three demand classes; its
internal axes are optimized for magnitude prediction, not class
separation. Second, 32 PCA dimensions let a tree ensemble learn
interactions across the embedding (for example, branching on one
component after another), whereas reducing to a single forecast scalar
hands the classifier a one-dimensional summary and removes any
possibility of such cross-component splits. Third, interpreting a
negative result from forecast-scalar-as-feature is ambiguous: a failure
could mean either the intermediate representation is uninformative, or
the representation is fine but the forecast head channels it poorly for
classification. Working directly on the intermediate representation
separates these two questions. A forecast-head-as-feature comparison is
a reasonable follow-up experiment.

We emphasize that this is a null result \emph{for this specific use case},
not a critique of TimesFM as a forecasting model~\cite{das2024timesfm}. The
model may well be the right tool for other tasks, particularly univariate
forecasting in domains with less hand-crafted feature engineering. Our
finding is specific to using its penultimate-layer representations as
tabular features on structured transit event data where short-range history
is already the dominant signal.

\subsection{Generalization of the forward-fill lag trap}

The target leakage pattern we document is not specific to bus ridership. It
arises whenever three conditions coincide: (1) the underlying signal is
sparse event data, (2) forward-fill imputation produces a dense series, and
(3) lag features are computed on the dense series with row offsets
smaller than the typical inter-event gap. All three conditions are
satisfied in many real-world domains:

\begin{itemize}
  \item \textbf{Medical vital signs.} Checkups produce sparse blood pressure
    or glucose readings that are frequently forward-filled onto dashboards
    at 1-minute resolution. A 5-minute lag feature on the dense series
    equals the current value for the vast majority of rows between two
    checkups.
  \item \textbf{Financial tick data.} Quote updates arrive at irregular
    intervals; forward-filling onto a 1-second grid and computing lag
    features exhibits the same pattern between ticks.
  \item \textbf{IoT sensor streams.} Sensors that report on thresholds
    (door open, motion detected, temperature exceeding a setpoint) produce
    sparse events often carried forward for aggregate reporting and
    downstream analytics.
  \item \textbf{Event logs in general.} Any time a log of discrete events
    is upsampled to a regular grid to match other streams, the resulting
    dense series is vulnerable.
\end{itemize}

In all these cases, the fix we propose (compute features at the natural
event granularity and forward-fill the feature values, not the raw series
before feature computation) generalizes directly. The key design pattern
is to \emph{preserve the original event marker} (our
\texttt{is\_stop\_event} flag) through the preprocessing pipeline so that
downstream feature engineering can recover the sparse structure after
imputation has filled in the gaps.

\subsection{Practical recommendations for transit operators}

For transit operators seeking to deploy a demand classification pipeline
similar to ours, we recommend:
\begin{enumerate}
  \item Compute lag features at the stop-event level, not the 1~Hz level.
    The computational overhead is minimal and the payoff (in avoided
    leakage and real predictive accuracy) is substantial.
  \item Use temporal evaluation splits (train on past missions, test on
    future missions) rather than random mission-level splits. The pooled
    evaluation overestimates deployment accuracy because the test
    distribution is identical to train.
  \item Start with Random Forest as the production baseline before
    reaching for larger models. We find RF and XGBoost perform nearly
    identically on this task, and both are easier to debug, interpret,
    and deploy than foundation models.
  \item Do not expect foundation model embeddings to help ``for free''
    as drop-in features. They may require fine-tuning, longer context
    windows, or task-specific adapter layers to contribute beyond
    well-designed hand-crafted features.
\end{enumerate}

\subsection{Limitations}

Our study has several limitations. First, we evaluate on a single city
(Zurich) and a single transit mode (trolleybus). The stop-level lag pattern
and the leakage fix should generalize, but the specific F1 numbers will
not. Second, we do not fine-tune TimesFM on ZTBus data; our null result
applies to zero-shot feature extraction, not to fine-tuned use. Third, our
classification target is tercile-binned passenger count, which discards
information about the exact count magnitude. A regression formulation
might yield different conclusions, particularly for the high-demand tail
where the absolute count matters more than which bucket it falls in.
Fourth, we do not deploy the model in a real-time operational setting;
the forecast horizon in production deployment would need to be explicitly
modeled (here we classify the current row rather than forecast a future
row). Fifth, we use mean-pooling to aggregate TimesFM embeddings, which is
simple but lossy; a more sophisticated aggregation strategy may yield
different results on the foundation model side.

\textbf{Tree regularization.} We train Decision Tree and Random Forest
with scikit-learn defaults: unlimited depth (\texttt{max\_depth=None}),
no pruning, and no explicit minimum-leaf-size constraint. XGBoost uses
\texttt{max\_depth=6}, a standard library default. Temporal
out-of-sample performance (macro-F1 = 0.8287 for RF on 2022 held-out
missions) suggests the learned structure is not trivially overfit to
the pre-2022 training partition, but we do not verify generalization
beyond ZTBus. An unbounded RF may encode splits idiosyncratic to the
Zurich trolleybus operating envelope (routes, fleet, weather, COVID-era
ridership). Regularized variants (bounded depth, minimum-leaf-samples,
cost-complexity pruning) and evaluation on a disjoint transit agency
(a different city or bus system entirely) are natural follow-ups for
operators considering cross-system deployment.

\textbf{Request-stop conventions.} ZTBus records trolleybus missions in
Zurich, where scheduled stops are serviced by default in normal
operation: the on-board ITCS reports a passenger-count update at each
stop event when the doors open. Our stop-level lag features implicitly
assume this default-service convention. In transit systems with
request-stop conventions, where drivers only stop when a passenger
signals demand (for example the MBTA Silver Line and many suburban
routes), the event sequence becomes endogenous to demand itself: stops
appear preferentially where passengers are, inducing selection bias
into any feature computed at stop-event granularity. The forward-fill
lag trap fix we propose would still eliminate the leakage artifact, but
the resulting features would carry a demand-selection confound that is
absent in a default-service regime. This is a constraint on
cross-system generalization, not on the ZTBus result.
